\chapter[Metody wariacyjne i twierdzenie Laxa--Milgrama]{Metody wariacyjne: słabe rozwiązania, twierdzenie Laxa--Milgrama i minimalizacja energii}
\label{chap:pde-variational}

Równania eliptyczne rzadko dają się rozwiązać za pomocą jawnych wzorów.
Wcześniej nauczyliśmy się jednak kontrolować słabe pochodne,
operator śladu i normy przestrzeni Sobolewa. Teraz pokażemy, że
te narzędzia pozwalają udowadniać istnienie oraz jednoznaczność
rozwiązań bez konstruowania ich punkt po punkcie.

Kluczowe będą dwa podejścia. Pierwsze polega na rozwiązywaniu
równania liniowego w przestrzeni Hilberta; jego podstawą jest
twierdzenie Laxa--Milgrama. Drugie polega na minimalizacji energii,
czasem również nieliniowej. Będziemy wyraźnie odróżniać
\emph{formę dwuliniową} od \emph{funkcjonału energii}: jeśli forma nie
jest symetryczna, to zwykły funkcjonał kwadratowy nie daje
na ogół wyjściowego równania.

W całym rozdziale przestrzenie Hilberta oraz formy dwuliniowe są
rzeczywiste. Całkujemy względem miary Lebesgue'a
$\lambda_d$ na $\Omega\subset\R^d$, zaś na brzegu Lipschitzowskim
względem $\mathcal H^{d-1}$. Zapisy $\int_\Omega f\,\mathrm d\lambda_d(x)$
i $\int_{\partial\Omega}g\,\mathrm d\mathcal H^{d-1}(x)$
określają za każdym razem także miarę całkowania.

\section{Przestrzeń energetyczna i sformułowanie słabe}
\label{sec:pde11-weak}

Niech $\Omega\subset\R^d$ będzie obszarem ograniczonym Lipschitzowskim. Kładziemy $V=H_0^1(\Omega)=W_0^{1,2}(\Omega)$.
Ustalmy stałą $C_P>0$ taką, że
$\|v\|_2\le C_P\|\nabla v\|_2$ dla każdego $v\in V$
(wniosek~\ref{cor:pde10-Poincare}). Z tej nierówności wynika, że
\[ (u,v)_V=\int_\Omega\nabla u\cdot\nabla v\,\mathrm d\lambda_d(x) \]
wyznacza iloczyn skalarny: jest symetryczna i dwuliniowa,
$(u,u)_V=\|\nabla u\|_2^2\ge0$, a gdy $(u,u)_V=0$,
nierówność Poincarégo daje $\|u\|_2=0$, czyli $u=0$ w $V$.
Norma $\|u\|_V=\sqrt{(u,u)_V}$ jest równoważna normie
$H^1(\Omega)$. W istocie $\|u\|_{H^1}^2=\|u\|_2^2+
\|\nabla u\|_2^2\le(1+C_P^2)\|\nabla u\|_2^2$,
a $\|u\|_V\le\|u\|_{H^1}$ wynika z definicji normy.
Przestrzeń $H_0^1$ jest domknięciem liniowej przestrzeni
$C_c^\infty(\Omega)$ w $H^1$, a więc domkniętą podprzestrzenią
przestrzeni Hilberta $H^1$ (twierdzenie~\ref{thm:pde9-banach}).
Ciąg Cauchy'ego w normie $\|\cdot\|_V$ jest dzięki pierwszemu
oszacowaniu ciągiem Cauchy'ego w $H^1$, ma tam granicę w $V$
i zbiega do niej także w normie $\|\cdot\|_V$.
Zatem $(V,(\cdot,\cdot)_V)$ jest przestrzenią Hilberta.

\begin{definicja}[Przestrzeń dualna i słabe rozwiązanie]
\label{def:pde11-dual}
Niech $V$ będzie rzeczywistą przestrzenią Hilberta.
Jej \emph{przestrzenią dualną} $V^*$ nazywamy przestrzeń wszystkich
ciągłych funkcjonałów liniowych $F:V\to\R$, wyposażoną w normę
$\|F\|_{V^*}=\sup_{\|v\|_V\le1}|F(v)|$.
Ciągłość liniowego funkcjonału jest równoważna skończoności
tego supremum: z ciągłości w zerze istnieje $\delta>0$,
dla którego $\|v\|_V\le\delta$ implikuje $|F(v)|\le1$,
a po przeskalowaniu otrzymujemy $\|F\|_{V^*}\le1/\delta$.
W drugą stronę dla $v\ne0$ stosujemy definicję do
$v/\|v\|_V$, uzyskując $|F(v)|\le\|F\|_{V^*}\|v\|_V$;
dla $v=0$ nierówność jest równością.
Jednorodność i nierówność trójkąta wynikają z własności
modułu i supremum, a norma zerowa oznacza $F=0$.
Jeżeli $a:V\times V\to\R$ jest formą dwuliniową i $F\in V^*$,
to \emph{słabym rozwiązaniem zagadnienia} jest element $u\in V$
spełniający $a(u,v)=F(v)$ dla każdego $v\in V$.
Dla $V=H_0^1(\Omega)$ oznaczamy też $V^*=H^{-1}(\Omega)$;
w tej notacji $H^{-1}(\Omega)$ oznacza przestrzeń dualną do
$H_0^1(\Omega)$, a nie przestrzeń wszystkich dystrybucji
z ujemną regularnością na całej $\R^d$.
W tym rozdziale normę $H^{-1}$ liczymy względem normy
energetycznej $\|\cdot\|_V$. Jeśli $\|F\|_{(H^1)^*_{0}}$
oznacza normę funkcjonału na $H_0^1$ liczoną względem
pełnej normy $H^1$, to z zawierania odpowiednich kul
jednostkowych i ich skalowania mamy
\[
 \|F\|_{(H^1)^*_{0}}\le\|F\|_{V^*}
 \le\sqrt{1+C_P^2}\,\|F\|_{(H^1)^*_{0}}.
\]
Obie normy określają te same ciągłe funkcjonały.
Każdy $F\in H^{-1}(\Omega)$ określa dystrybucję przez
ograniczenie do $C_c^\infty(\Omega)$: dla testu o nośniku
w ustalonym zwartym $K\subset\Omega$ zachodzi
$|F(\varphi)|\le\|F\|_{V^*}\sqrt{d\lambda_d(K)}
\max_i\|\partial_i\varphi\|_\infty$.
Jest to oszacowanie ciągłości w topologii testów.
Ograniczenie jest iniektywne z gęstości testów w $V$.
\end{definicja}

\begin{przyklad}[Od równania Poissona do formy słabej]
\label{ex:pde11-poisson-weak}
Załóżmy chwilowo, że $u\in C^2(\overline\Omega)$,
$u|_{\partial\Omega}=0$, $f\in L^2(\Omega)$ oraz
$-\Delta u=f$ prawie wszędzie. Przez $C^2(\overline\Omega)$
rozumiemy tu funkcje klasy $C^2$ w otoczeniu domknięcia.
Funkcja $u$ i jej pochodne są ograniczone na zwartym
$\overline\Omega$, stąd $u\in H^1(\Omega)$.
Jej ślad jest równy restrykcji klasycznej: wybieramy gładkie
odcięcie równe jeden na otoczeniu $\overline\Omega$,
o zwartym nośniku w otoczeniu, gdzie $u$ jest klasy $C^2$.
Iloczyn przedłużamy przez zero na $\R^d$ i wygładzamy
splotem; przybliżenia zbiegają wraz z pierwszymi pochodnymi
jednostajnie na $\overline\Omega$ dzięki jednostajnej ciągłości.
Z ciągłości śladu w $H^1$ i skończonej miary brzegu
(twierdzenie~\ref{thm:pde10-trace}) ślad $u$ jest zatem zerowy.
Twierdzenie~\ref{thm:pde10-trace-zero} daje $u\in V$.
Dla $v\in C_c^\infty(\Omega)$ całkujemy przez części
w każdej współrzędnej (bez składników brzegowych z uwagi
na zwarty nośnik testu) i sumujemy. Otrzymujemy
$\int_\Omega\nabla u\cdot\nabla v\,\mathrm d\lambda_d(x)
=\int_\Omega fv\,\mathrm d\lambda_d(x)$.
Lewa strona jest ciągła względem $v$ w normie $H_0^1$,
ponieważ ogranicza ją $\|\nabla u\|_2\|\nabla v\|_2$.
Prawa strona również jest ciągła:
$|\int_\Omega fv\,\mathrm d\lambda_d|
\le C_P\|f\|_2\|\nabla v\|_2$.
Dla dowolnego $v\in V$ wybieramy $v_n\in C_c^\infty(\Omega)$
zbieżne w $H^1$, a więc także w normie $V$. Różnice
całek dla $v_n$ i $v$ są ograniczone odpowiednio przez
$\|\nabla u\|_2\|v_n-v\|_V$ i $C_P\|f\|_2\|v_n-v\|_V$.
Obie dążą do zera, co daje tożsamość dla wszystkich $v\in V$.
\end{przyklad}

W odwrotną stronę słaba tożsamość daje równość
$-\Delta u=f$ w sensie dystrybucji. Istotnie, dla
$v\in C_c^\infty(\Omega)$ z definicji pochodnej słabej
i pochodnej dystrybucji mamy
\[
 \langle-\Delta u,v\rangle
 =-\sum_{i=1}^d\langle\partial_i^2u,v\rangle
 =\sum_{i=1}^d\int_\Omega\partial_i u\,\partial_i v
       \,\mathrm d\lambda_d(x)
 =\int_\Omega fv\,\mathrm d\lambda_d(x).
\]
Dla ogólnego $F\in V^*$ ostatnią całkę zastępujemy
przez $F(v)$, czyli działanie dystrybucji zdefiniowanej
przez ograniczenie $F$ do testów.
Warunek $u|_{\partial\Omega}=0$ jest zakodowany w wyborze $V$;
na obszarze Lipschitzowskim jest równoważny
$\operatorname{Tr}u=0$ (twierdzenie~\ref{thm:pde10-trace-zero}).
Nie wynika stąd jeszcze, że $u\in H^2$ lub $u\in C^2$.
\emph{Istnienie słabe i regularność klasyczna to dwa osobne twierdzenia.}

\section{Dwa lematy o przestrzeniach Hilberta}
\label{sec:pde11-hilbert}

Najpierw udowodnimy fakty, od których zależy twierdzenie Laxa--Milgrama.
Nie zakładamy z góry, że każdy funkcjonał można zapisać
w postaci iloczynu skalarnego.

\begin{lemat}[Rzut na domknięty zbiór wypukły]
\label{lem:pde11-projection}
Niech $H$ będzie rzeczywistą przestrzenią Hilberta,
$K\subset H$ zbiorem niepustym, domkniętym i wypukłym,
a $x\in H$. Istnieje dokładnie jeden $y\in K$ taki, że
$\|x-y\|=\inf_{z\in K}\|x-z\|$. Ponadto
\[ (x-y,z-y)_H\le0\qquad\text{dla każdego }z\in K. \]
\end{lemat}
\begin{proof}
Oznaczmy $r=\inf_{z\in K}\|x-z\|$. Ponieważ $K$ jest
niepusty, $0\le r<\infty$. Z definicji infimum wybieramy
$y_n\in K$ z $r\le\|x-y_n\|<r+1/n$. Ponieważ $K$ jest wypukły,
$(y_n+y_m)/2\in K$. Tożsamość równoległoboku daje
\begin{align*}
\|y_n-y_m\|^2
&=2\|x-y_n\|^2+2\|x-y_m\|^2
 -4\left\|x-\frac{y_n+y_m}{2}\right\|^2\\
&\le 2\|x-y_n\|^2+2\|x-y_m\|^2-4r^2\longrightarrow0.
\end{align*}
Zatem $(y_n)$ jest ciągiem Cauchy'ego. Zupełność i domkniętość
zapewniają istnienie granicy $y\in K$; ciągłość normy
daje $\|x-y\|=r$.
Gdyby $y'\in K$ było drugim minimizatorem, ten sam rachunek
z $y_n=y$, $y_m=y'$ dałby $\|y-y'\|=0$.

Dla dowolnego $z\in K$ i $0\le t\le1$ mamy
$y+t(z-y)\in K$. Z minimalności $y$ wynika
$\|x-y-t(z-y)\|^2\ge\|x-y\|^2$.
Rozwijając kwadrat, dzieląc przez $t>0$ i przechodząc
do granicy $t\downarrow0$, otrzymujemy
$-2(x-y,z-y)_H\ge0$, czyli tezę.
Odwrotnie, jeżeli $y\in K$ spełnia tę nierówność
rzutową dla wszystkich $z\in K$, to
\[
 \|x-z\|^2=\|x-y\|^2-2(x-y,z-y)_H+\|z-y\|^2
 \ge\|x-y\|^2.
\]
Zatem nierówność charakteryzuje jedyny punkt minimalizujący.
Gdy $K$ jest podprzestrzenią liniową, podstawienie
$z=y+k$ i $z=y-k$ dla $k\in K$ daje $(x-y,k)_H=0$.
\end{proof}

\begin{twierdzenie}[Twierdzenie Riesza o reprezentacji]
\label{thm:pde11-riesz}
Niech $H$ będzie rzeczywistą przestrzenią Hilberta.
Dla każdego $F\in H^*$ istnieje dokładnie jeden $w_F\in H$ taki, że
$F(v)=(w_F,v)_H$ dla wszystkich $v\in H$.
Ponadto $\|w_F\|_H=\|F\|_{H^*}$.
\end{twierdzenie}
\begin{proof}
Dla $F=0$ wektor $w_F=0$ spełnia żądaną tożsamość
i równość norm; jeśli inny wektor ją spełnia, podstawienie
tego wektora za $v$ daje jego zerową normę.
Jeżeli $F\ne0$, to $K=\ker F$ jest podprzestrzenią
liniową, domkniętą jako przeciwobraz $\{0\}$ przez
ciągłe $F$ i właściwą, ponieważ $F$ nie jest zerowy. Wybierzmy $x\notin K$.
Z lematu~\ref{lem:pde11-projection} otrzymujemy rzut
$y\in K$ punktu $x$ i kładziemy $e=x-y\ne0$.
Z charakterystyki rzutu, podstawiając $z=y\pm k$,
$k\in K$, dostajemy $(e,k)_H=0$.
Dla dowolnego $v\in H$ element
$k=v-[F(v)/F(e)]e$ należy do $K$, ponieważ
$F(e)=F(x)\ne0$. Stąd
$0=(e,k)_H=(e,v)_H-[F(v)/F(e)]\|e\|_H^2$.
Otrzymujemy $F(v)=(w_F,v)_H$ dla
$w_F=F(e)e/\|e\|_H^2$.
Jeżeli $w_1,w_2$ reprezentują $F$, to
$(w_1-w_2,v)_H=0$ dla każdego $v\in H$; podstawiając
$v=w_1-w_2$, dostajemy $\|w_1-w_2\|^2=0$.
Nierówność Cauchy'ego--Schwarza daje
$|F(v)|\le\|w_F\|_H$ dla $\|v\|_H\le1$, więc
$\|F\|_{H^*}\le\|w_F\|_H$; biorąc
$v=w_F/\|w_F\|$ w przypadku $w_F\ne0$,
uzyskujemy $|F(v)|=\|w_F\|_H$, czyli nierówność odwrotną.
Jednoznaczność zapewnia, że przyporządkowanie
$F\mapsto w_F$ jest liniowe: kombinacja liniowa
reprezentantów reprezentuje tę samą kombinację funkcjonałów.
Jest ono izometrią na $H$, ponieważ każdy $w\in H$
reprezentuje ciągły funkcjonał $v\mapsto(w,v)_H$.
\end{proof}

\section{Twierdzenie Laxa--Milgrama z pełnym dowodem}
\label{sec:pde11-lax}

\begin{definicja}[Ograniczoność i koercywność]
\label{def:pde11-coercivity}
Niech $H$ będzie rzeczywistą przestrzenią Hilberta.
Formę dwuliniową $a:H\times H\to\R$ nazywamy
\emph{ograniczoną}, jeżeli istnieje $M\ge0$ takie, że
$|a(u,v)|\le M\|u\|_H\|v\|_H$ dla wszystkich $u,v\in H$.
Nazywamy ją \emph{koercywną}, jeżeli istnieje $\alpha>0$,
dla którego $a(v,v)\ge\alpha\|v\|_H^2$ dla każdego $v\in H$.
Koercywność nie jest tym samym co symetria: forma
może być koercywna i jednocześnie niesymetryczna.
Dla przykładu na $\R^2$ forma
$a(u,v)=u\cdot v+u_1v_2-u_2v_1$ spełnia
$a(v,v)=|v|^2$, ale $a(e_1,e_2)=1$ i $a(e_2,e_1)=-1$.
Jest ograniczona, gdyż
$|u_1v_2-u_2v_1|\le |u|\,|v|$ z nierówności
Cauchy'ego--Schwarza dla $(u_1,u_2)$ i $(v_2,-v_1)$.
Można więc przyjąć $M=2$ i $\alpha=1$.
\end{definicja}

\begin{twierdzenie}[Twierdzenie Laxa--Milgrama]
\label{thm:pde11-LM}
Niech $H$ będzie rzeczywistą przestrzenią Hilberta,
$a:H\times H\to\R$ ograniczoną formą dwuliniową
ze stałą $M>0$, koercywną ze stałą $\alpha>0$,
a $F\in H^*$. Wtedy istnieje dokładnie jeden element $u\in H$
spełniający
\[ a(u,v)=F(v)\qquad\text{dla każdego }v\in H. \]
Ponadto $\|u\|_H\le\alpha^{-1}\|F\|_{H^*}$.
Nie wymaga się symetrii formy $a$.
\end{twierdzenie}
\begin{proof}
Jeśli $H=\{0\}$, jedynym rozwiązaniem jest $u=0$,
a jedyny funkcjonał to $F=0$; oszacowanie zachodzi.
Dalej zakładamy $H\ne\{0\}$. Dla dowolnego $v\ne0$
koercywność i ograniczoność dają
$\alpha\|v\|_H^2\le a(v,v)\le M\|v\|_H^2$,
a zatem $0<\alpha\le M$.

\textbf{Krok 1: operator reprezentujący formę.}
Dla ustalonego $u\in H$ odwzorowanie $v\mapsto a(u,v)$
jest funkcjonałem liniowym i ograniczonym o normie
co najwyżej $M\|u\|_H$. Z twierdzenia Riesza
(\ref{thm:pde11-riesz}) istnieje dokładnie jeden
$Tu\in H$ spełniający
$(Tu,v)_H=a(u,v)$ dla każdego $v\in H$.
Dla $u,w\in H$ i $s,r\in\R$ mamy
$(T(su+rw),v)_H=(sTu+rTw,v)_H$ dla każdego $v$,
z liniowości $a$ w pierwszej zmiennej.
Podstawiając za $v$ różnicę tych wektorów, otrzymujemy
$T(su+rw)=sTu+rTw$, czyli liniowość $T$.
Równość norm w twierdzeniu Riesza daje
$\|Tu\|_H\le M\|u\|_H$, a zatem $T$ jest ograniczony.
Koercywność oznacza $(Tu,u)_H\ge\alpha\|u\|_H^2$.
Z kolei $F(v)=(f,v)_H$ dla jednoznacznie wyznaczonego
$f\in H$ z $\|f\|_H=\|F\|_{H^*}$.
Równanie słabe jest równoważne $(Tu-f,v)_H=0$
dla każdego $v\in H$, czyli $Tu=f$ (bierzemy $v=Tu-f$).

\textbf{Krok 2: konstrukcja kontrakcji.}
Dla parametru $t>0$ zdefiniujmy
$S_t(u)=u-t(Tu-f)$. Dla $u,w\in H$, oznaczając $z=u-w$,
obliczamy dokładnie
\begin{align*}
\|S_t(u)-S_t(w)\|_H^2
&=\|z-tTz\|_H^2\\
&=\|z\|_H^2-2t(Tz,z)_H+t^2\|Tz\|_H^2\\
&\le (1-2t\alpha+t^2M^2)\|z\|_H^2.
\end{align*}
Wybierzmy $0<t<2\alpha/M^2$.
Z nierówności $\alpha\le M$ wynika
\[
 1-2t\alpha+t^2M^2
 =M^2(t-\alpha/M^2)^2+1-\alpha^2/M^2\ge0.
\]
Warunek $0<t<2\alpha/M^2$ daje zarazem
$t(tM^2-2\alpha)<0$, więc
$q=\sqrt{1-2t\alpha+t^2M^2}\in[0,1)$.
Po wyciągnięciu pierwiastka z oszacowania dostajemy
$\|S_t(u)-S_t(w)\|_H\le q\|u-w\|_H$,
czyli $S_t$ jest kontrakcją.
Najmniejszą wartość tego oszacowania stałej kontrakcji
otrzymujemy dla $t=\alpha/M^2$, mianowicie
$q=\sqrt{1-\alpha^2/M^2}$; nie musi to być rzeczywista
najmniejsza stała Lipschitza odwzorowania $S_t$.

\textbf{Krok 3: punkt stały bez pomijania granicy.}
Ustalmy dowolne $u_0\in H$ i zdefiniujmy
$u_{n+1}=S_t(u_n)$.
Jeśli $q=0$, odwzorowanie $S_t$ jest stałe,
zatem $u_2=S_t(u_1)=S_t(u_0)=u_1$ i $u_1$ jest
punktem stałym. Niech teraz $0<q<1$.
Z kontrakcji, przez indukcję względem $n$, mamy
$\|u_{n+1}-u_n\|_H\le q^n\|u_1-u_0\|_H$:
dla $n=0$ jest równość, a krok indukcyjny wynika
z zastosowania $S_t$ do $u_{n+1}$ i $u_n$.
Dla $m>n$ sumujemy te oszacowania:
\[
 \|u_m-u_n\|_H\le
 \sum_{j=n}^{m-1}\|u_{j+1}-u_j\|_H
 \le\frac{q^n}{1-q}\|u_1-u_0\|_H.
\]
Prawa strona dąży do zera niezależnie od $m>n$.
Zupełność $H$ zapewnia istnienie $u=\lim_nu_n$.
Ponieważ $S_t$ jest ciągłe, $S_t(u_n)\to S_t(u)$,
a jednocześnie $S_t(u_n)=u_{n+1}\to u$.
Stąd $u=S_t(u)$, a ponieważ $t>0$, mamy $Tu=f$.
Przejście z $m$ do nieskończoności daje dodatkowo
oszacowanie błędu
$\|u-u_n\|_H\le q^n(1-q)^{-1}\|u_1-u_0\|_H$.
Jeżeli $Tu=f=Tw$, to z koercywności
$\alpha\|u-w\|_H^2\le(T(u-w),u-w)_H=0$;
stąd $u=w$.

\textbf{Krok 4: oszacowanie rozwiązania.}
Podstawiając $v=u$ w tożsamości słabej, otrzymujemy
$\alpha\|u\|_H^2\le a(u,u)=F(u)
\le\|F\|_{H^*}\|u\|_H$.
Dla $u\ne0$ dzielimy przez dodatnią liczbę $\alpha\|u\|_H$;
dla $u=0$ lewa strona żądanego oszacowania jest zerowa.
Ponadto $T$ jest bijekcją: dla dowolnego $f\in H$
stosujemy udowodnione twierdzenie do funkcjonału
$v\mapsto(f,v)_H$, a jednoznaczność daje iniektywność.
Z równości norm w twierdzeniu Riesza
$\|T^{-1}f\|_H\le\alpha^{-1}\|f\|_H$.
Odwzorowanie odwrotne jest liniowe jako odwrotność
bijekcji liniowej i jest ograniczone.
\end{proof}

\begin{wniosek}[Stabilność względem prawej strony]
\label{cor:pde11-stability}
Przy założeniach twierdzenia~\ref{thm:pde11-LM},
niech $u_i\in H$ będą rozwiązaniami odpowiadającymi
$F_i\in H^*$, $i=1,2$. Wówczas
$\|u_1-u_2\|_H\le\alpha^{-1}\|F_1-F_2\|_{H^*}$.
\end{wniosek}
\begin{proof}
Odejmujemy słabe równania: $a(u_1-u_2,v)=(F_1-F_2)(v)$.
Dla $w=u_1-u_2$ bierzemy $v=w$, otrzymując
$\alpha\|w\|_H^2\le a(w,w)=(F_1-F_2)(w)
\le\|F_1-F_2\|_{H^*}\|w\|_H$.
Jeśli $w\ne0$, dzielimy przez $\alpha\|w\|_H$;
w przeciwnym przypadku oszacowanie zachodzi.
Jest to ilościowa stabilność względem danych:
wraz z istnieniem i jednoznacznością oznacza poprawność
zagadnienia w sensie Hadamarda dla ustalonej formy $a$.
\end{proof}

\begin{uwaga}[Metoda iteracyjna a dowód istnienia]
Punkt stały w dowodzie dostarcza jednocześnie algorytmu:
$u_{n+1}=u_n-t(Tu_n-f)$, o zbieżności geometrycznej
przy podanym wyborze $t$. Sam operator $T$ zależy jednak
od wybranego iloczynu skalarnego w $H$, toteż bez
skończeniewymiarowej reprezentacji nie jest to automatycznie
praktyczna metoda numeryczna.
\end{uwaga}

\section{Równanie eliptyczne z mierzalnymi współczynnikami}
\label{sec:pde11-elliptic}

\begin{twierdzenie}[Zagadnienie Dirichleta w postaci dywergencyjnej]
\label{thm:pde11-elliptic}
Niech $\Omega\subset\R^d$ będzie obszarem ograniczonym Lipschitzowskim. Załóżmy, że $A:\Omega\to\R^{d\times d}$
ma mierzalne, istotnie ograniczone współczynniki,
a $c\in L^\infty(\Omega)$ jest nieujemna prawie wszędzie.
Dla macierzy używamy normy operatorowej względem
normy euklidesowej, a dla wektorów normy euklidesowej.
Ich normy $L^\infty$ oznaczają istotne suprema
tych norm punktowych.
Załóżmy ponadto, że dla pewnego $\nu>0$ oraz
prawie wszystkich $x\in\Omega$ zachodzi
$\xi^TA(x)\xi\ge\nu|\xi|^2$ dla każdego $\xi\in\R^d$.
Wówczas dla każdego $F\in H^{-1}(\Omega)$ istnieje
jednoznaczne $u\in H_0^1(\Omega)$ spełniające
\[
\int_\Omega\big(A\nabla u\cdot\nabla v+cuv\big)
\,\mathrm d\lambda_d(x)=F(v)
\quad\text{dla wszystkich }v\in H_0^1(\Omega).
\]
Przy normie $\|v\|_V=\|\nabla v\|_2$ zachodzi
$\|u\|_V\le\nu^{-1}\|F\|_{V^*}$.
Nie zakłada się symetrii macierzy $A(x)$.
\end{twierdzenie}
\begin{proof}
Na $V=H_0^1(\Omega)$ zdefiniujmy
$a(u,v)=\int_\Omega(A\nabla u\cdot\nabla v+cuv)
\,\mathrm d\lambda_d(x)$.
Oznaczmy $K=\operatorname*{ess\,sup}_{x\in\Omega}
\|A(x)\|_{\mathrm{op}}$ oraz $C=\|c\|_\infty$.
Norma macierzy jest funkcją ciągłą jej współczynników,
więc $x\mapsto\|A(x)\|_{\mathrm{op}}$ jest mierzalna.
Ponadto $\|A(x)\|_{\mathrm{op}}\le
(\sum_{i,j}|A_{ij}(x)|^2)^{1/2}$, z nierówności
Cauchy'ego--Schwarza dla każdego wiersza i sumowania
kwadratów. Zatem $K<\infty$.
Poza jednym zbiorem miary zero mamy
$|A(x)\xi|\le K|\xi|$ dla wszystkich $\xi\in\R^d$.
Składniki podcałkowe są mierzalne, niezależne od
reprezentantów i całkowalne, ponieważ ich moduły
ograniczają odpowiednio $K|\nabla u||\nabla v|$
oraz $C|u||v|$, należące do $L^1$ z Höldera.
Liniowość całki zapewnia dwuliniowość formy.
Z Cauchy'ego--Schwarza i nierówności Poincar\'ego
otrzymujemy
\begin{align*}
|a(u,v)|
&\le K\|\nabla u\|_2\|\nabla v\|_2
+C\|u\|_2\|v\|_2\\
&\le(K+CC_P^2)\|u\|_V\|v\|_V.
\end{align*}
Tak więc $a$ jest formą ograniczoną.
Ponieważ $c\ge0$, eliptyczność zapewnia
$a(v,v)\ge\nu\int_\Omega|\nabla v|^2
\,\mathrm d\lambda_d=\nu\|v\|_V^2$.
Wszystkie założenia twierdzenia~\ref{thm:pde11-LM} są spełnione,
więc otrzymujemy istnienie, jednoznaczność i oszacowanie.
\end{proof}

\begin{uwaga}[Znaczenie sformułowania dystrybucyjnego]
Dla $F_f(v)=\int_\Omega fv\,\mathrm d\lambda_d(x)$,
$f\in L^2(\Omega)$, zachodzi
$\|F_f\|_{V^*}\le C_P\|f\|_2$.
W tym przypadku rozwiązanie spełnia w $\mathcal D'(\Omega)$
$-\operatorname{div}(A\nabla u)+cu=f$.
Dokładniej, dla $\varphi\in C_c^\infty(\Omega)$
definicja pochodnej dystrybucyjnej daje
\[
 \langle-\operatorname{div}(A\nabla u)+cu,\varphi\rangle
 =\int_\Omega(A\nabla u\cdot\nabla\varphi+cu\varphi)
       \,\mathrm d\lambda_d(x)
 =F_f(\varphi).
\]
Dla ogólnego $F\in V^*$ ostatnia wartość to $F(\varphi)$,
działanie dystrybucji z definicji~\ref{def:pde11-dual}.
Odwrotnie, jeśli $u\in V$ spełnia tę równość
dystrybucyjnie, to gęstość testów w $V$ oraz
ograniczoność formy i funkcjonału rozszerzają ją na $V$.
Nie ma tu obowiązku definiowania klasycznej dywergencji
$A\nabla u$: przy mierzalnych współczynnikach produkt
$A\nabla u$ może nie mieć słabych pochodnych w $L^2$,
ale ma sens jako wektor z $L^2$, którego dywergencja
jest dystrybucją.
\end{uwaga}

\begin{przyklad}[Jawne rozwiązanie w jednym wymiarze]
\label{ex:pde11-1d}
Dla $\Omega=(0,1)$, $A=1$, $c=0$ i $f=1$
szukamy $u\in H_0^1(0,1)$ spełniającego
$\int_0^1u'(x)v'(x)\,\mathrm d\lambda_1(x)
=\int_0^1v(x)\,\mathrm d\lambda_1(x)$.
Dla $u(x)=x(1-x)/2$ mamy $u'(x)=1/2-x$,
$u''(x)=-1$ i $u(0)=u(1)=0$.
Funkcja jest w $H^1(0,1)$, a z jednowymiarowego
oszacowania śladu i twierdzenia~\ref{thm:pde10-trace-zero}
należy do $H_0^1(0,1)$.
Dla testów $v$ całkowanie przez części daje
$\int_0^1 u'v'\,\mathrm d\lambda_1
=-\int_0^1 u''v\,\mathrm d\lambda_1
=\int_0^1 v\,\mathrm d\lambda_1$.
Przejście przez gęstość jak w przykładzie
~\ref{ex:pde11-poisson-weak} daje tożsamość dla wszystkich $v\in V$. Twierdzenie~\ref{thm:pde11-elliptic}
wykazuje, że nie istnieje żadne inne rozwiązanie w $H_0^1$.
\end{przyklad}

\begin{twierdzenie}[Koercywność przy wyrazach niższego rzędu]
\label{thm:pde11-lower-order}
Niech $\Omega\subset\R^d$ będzie ograniczonym obszarem
Lipschitzowskim, a $V=H_0^1(\Omega)$ z normą
$\|u\|_V=\|\nabla u\|_{L^2(\Omega)}$. Niech
$A\in L^\infty(\Omega;\R^{d\times d})$ spełnia
$\xi^TA(x)\xi\ge\nu|\xi|^2$ dla pewnego $\nu>0$,
wszystkich $\xi\in\R^d$ i prawie wszystkich $x\in\Omega$.
Niech $b\in L^\infty(\Omega;\R^d)$ i
$c\in L^\infty(\Omega;\R)$; oznaczmy
$c_-(x)=\max\{-c(x),0\}$ i przez $C_P$ dowolną
stałą w nierówności $\|u\|_2\le C_P\|\nabla u\|_2$
dla $u\in V$. Jeśli
\[
\alpha:=\nu-C_P\|b\|_\infty-C_P^2\|c_-\|_\infty>0,
\]
to forma
\[
 a(u,v)=\int_\Omega\bigl(A\nabla u\cdot\nabla v
 +(b\cdot\nabla u)v+cuv\bigr)\,\mathrm d\lambda_d(x)
\]
jest ograniczona i koercywna na $V$ ze stałą co
najmniej $\alpha$. Dla każdego $F\in V^*$ równanie
$a(u,v)=F(v)$, $v\in V$, ma dokładnie jedno rozwiązanie
i zachodzi $\|u\|_V\le\alpha^{-1}\|F\|_{V^*}$.
\end{twierdzenie}
\begin{proof}
Przyjmujemy tę samą konwencję norm macierzy i wektorów
co w twierdzeniu~\ref{thm:pde11-elliptic}.
Eliptyczność rozumiemy poza wspólnym zbiorem miary zero;
jeśli nierówność przyjęto osobno dla każdego $\xi$,
wystarczy najpierw przeciąć zbiory pełnej miary
dla $\xi\in\mathbb Q^d$, a potem użyć ciągłości
funkcji kwadratowej względem $\xi$.
Z nierówności Cauchy'ego--Schwarza i Höldera mamy
\[
 \begin{aligned}
 |a(u,v)|&\le\|A\|_\infty\|\nabla u\|_2\|\nabla v\|_2\\
 &\quad+\|b\|_\infty\|\nabla u\|_2\|v\|_2
       +\|c\|_\infty\|u\|_2\|v\|_2\\
 &\le(\|A\|_\infty+C_P\|b\|_\infty
       +C_P^2\|c\|_\infty)\|u\|_V\|v\|_V,
 \end{aligned}
\]
gdzie ostatni krok to nierówność Poincarégo.
Oszacowanie uzasadnia również całkowalność każdego
składnika i ograniczoność formy.
Z kolei dla $u\in V$ dodatnia część $c$ nie może
obniżać wartości $a(u,u)$, a pozostałe wyrazy dają
\begin{align*}
 a(u,u)&\ge\nu\|\nabla u\|_2^2
 -\|b\|_\infty\|\nabla u\|_2\|u\|_2
 -\|c_-\|_\infty\|u\|_2^2\\
 &\ge\alpha\|u\|_V^2.
\end{align*}
Twierdzenie Laxa--Milgrama~\ref{thm:pde11-LM}
kończy dowód. Warunek ma charakter wystarczający;
ujemna część potencjału lub dryft mogą być większe,
jeżeli znamy dodatkową strukturę formy.
\end{proof}

\begin{przyklad}[Dryft bez ograniczenia małości]
\label{ex:pde11-divfree-drift}
Niech $b\in W^{1,\infty}(\Omega;\R^d)$ spełnia
$\operatorname{div}b=0$ prawie wszędzie oraz
$A=I$, $c=0$. Dla każdego $u\in C_c^\infty(\Omega)$
całkowanie przez części daje
$\int_\Omega(b\cdot\nabla u)u\,\mathrm d\lambda_d
=\frac12\int_\Omega b\cdot\nabla(u^2)\,\mathrm d\lambda_d=0$.
Istotnie, $u^2$ jest testem i z definicji
dywergencji dystrybucyjnej
$\int_\Omega b\cdot\nabla(u^2)\,\mathrm d\lambda_d
=-\langle\operatorname{div}b,u^2\rangle=0$.
Dla dowolnego $u\in H_0^1(\Omega)$ wybieramy
$u_n\in C_c^\infty(\Omega)$ zbieżne do $u$ w $H^1$.
Z nierówności Höldera
\[
 \begin{aligned}
 &\left|\int_\Omega
 [(b\cdot\nabla u_n)u_n-(b\cdot\nabla u)u]
               \,\mathrm d\lambda_d\right|\\
 &\qquad\le\|b\|_\infty
 (\|\nabla(u_n-u)\|_2\|u_n\|_2
   +\|\nabla u\|_2\|u_n-u\|_2)\longrightarrow0.
 \end{aligned}
\]
Ciąg norm $\|u_n\|_2$ jest ograniczony dzięki zbieżności.
Otrzymujemy tożsamość dla wszystkich $u\in V$. Dlatego
$a(u,u)=\|\nabla u\|_2^2$ i forma jest koercywna
ze stałą $1$ bez założenia, że $\|b\|_\infty$ jest mała.
Sam dryft nie jest więc przeszkodą: znaczenie ma również
jego struktura i warunki brzegowe.
\end{przyklad}

\section{Niezerowe dane brzegowe: metoda podnoszenia}
\label{sec:pde11-lifting}

\begin{twierdzenie}[Warunek Dirichleta z danymi brzegowymi]
\label{thm:pde11-lifting}
Niech $\Omega,A,c,\nu$ spełniają założenia
twierdzenia~\ref{thm:pde11-elliptic}. Niech
$f\in H^{-1}(\Omega)$ oraz $g\in H^1(\Omega)$.
Istnieje dokładnie jedna funkcja $u\in H^1(\Omega)$
taka, że $u-g\in H_0^1(\Omega)$ oraz
\[
\int_\Omega(A\nabla u\cdot\nabla v+cuv)
\,\mathrm d\lambda_d(x)=f(v)
\quad(v\in H_0^1(\Omega)).
\]
Jej ślad jest równy $\operatorname{Tr}g$.
Istnieje stała $C$, zależna wyłącznie od
$\Omega,\nu,\|A\|_\infty,\|c\|_\infty$, taka, że
$\|u\|_{H^1}\le C(\|f\|_{H^{-1}}+\|g\|_{H^1})$.
\end{twierdzenie}
\begin{proof}
Podstawiamy $u=w+g$. Żądany warunek brzegowy jest równoważny
$w\in V=H_0^1(\Omega)$, a równanie przyjmuje postać
$a(w,v)=G(v)$, gdzie $G(v)=f(v)-a(g,v)$.
Ograniczoność współczynników, Cauchy--Schwarz i
Poincar\'e dla $v\in V$ zapewniają
$|a(g,v)|\le C_0\|g\|_{H^1}\|v\|_V$.
Wobec tego $G\in V^*$ i
$\|G\|_{V^*}\le\|f\|_{V^*}+C_0\|g\|_{H^1}$.
Twierdzenie Laxa--Milgrama daje jednoznaczne $w\in V$ oraz
$\|w\|_V\le\nu^{-1}\|G\|_{V^*}$.
Ponieważ $\|w\|_{H^1}\le(1+C_P^2)^{1/2}\|w\|_V$,
uzyskujemy żądane oszacowanie dla $u=w+g$.
Równość śladów wynika z
$\operatorname{Tr}(u-g)=0$ i twierdzenia
\ref{thm:pde10-trace-zero}.
\end{proof}

\begin{uwaga}[Dane na samym brzegu]
Dane $\varphi$ na $\partial\Omega$ można podać niezależnie,
ale nie każde $\varphi\in L^2(\partial\Omega)$ jest śladem
funkcji z $H^1(\Omega)$. W twierdzeniu zakładamy explicite
istnienie podniesienia $g\in H^1(\Omega)$.
Przestrzeń danych brzegowych $H^{1/2}(\partial\Omega)$
z normą ilorazową i jej liniowe podnoszenie zdefiniowano
już w~\ref{def:pde10-halftrace} i
\ref{lem:pde10-halftrace-rightinverse}. Równoważność tej
normy z normą Gagliarda pozostaje oddzielnym twierdzeniem.
\end{uwaga}


\begin{uwaga}[Wersja z danymi zadanymi na brzegu]
Twierdzenie~\ref{thm:pde11-lifting} sformułowano
z wybranym podniesieniem $g\in H^1(\Omega)$.
Z lematu~\ref{lem:pde10-halftrace-rightinverse} można
bez nowych argumentów otrzymać dane
$\varphi\in H^{1/2}(\partial\Omega)$ i oszacowanie
przez ich normę śladową. Uogólnienie to, także
z wyrazami niższego rzędu, zostaje jawnie sformułowane
w twierdzeniu~\ref{thm:pde13-nonhom}; nie powielamy
tutaj jego dowodu.
\end{uwaga}

\section{Zagadnienie Neumanna i warunek zgodności}
\label{sec:pde11-neumann}

Dla laplasjanu z warunkiem Neumanna nie można oczekiwać
jednoznaczności: dodanie stałej do rozwiązania nie zmienia
gradientu ani pochodnej normalnej. Warunek całkowy
pojawia się natomiast już w sformułowaniu słabym.

\begin{twierdzenie}[Słabe zagadnienie Neumanna]
\label{thm:pde11-neumann}
Niech $\Omega\subset\R^d$ będzie obszarem ograniczonym Lipschitzowskim, $f\in L^2(\Omega)$ oraz
$h\in L^2(\partial\Omega)$.
Rozważmy równanie słabe
\[
\int_\Omega\nabla u\cdot\nabla v\,\mathrm d\lambda_d(x)
=\int_\Omega fv\,\mathrm d\lambda_d(x)
+\int_{\partial\Omega}h\operatorname{Tr}v
\,\mathrm d\mathcal H^{d-1}(x)
\quad(v\in H^1(\Omega)).
\]
Istnieje rozwiązanie $u\in H^1(\Omega)$ wtedy i tylko wtedy, gdy
\[
\int_\Omega f\,\mathrm d\lambda_d(x)
+\int_{\partial\Omega}h\,\mathrm d\mathcal H^{d-1}(x)=0.
\]
Gdy ono istnieje, wszystkie rozwiązania różnią się o stałą.
Po narzuceniu normalizacji $\int_\Omega u\,\mathrm d\lambda_d=0$
rozwiązanie jest jednoznaczne.
\end{twierdzenie}
\begin{proof}
\textbf{Konieczność.} Funkcja $v\equiv1$ należy do $H^1(\Omega)$
i ma zerowy gradient, więc jej podstawienie daje warunek całkowy.

\textbf{Wystarczalność.} Niech
$V_\circ=\{v\in H^1(\Omega):\int_\Omega v\,\mathrm d\lambda_d=0\}$.
Jest to domknięta podprzestrzeń $H^1$, więc przestrzeń Hilberta.
Nierówność Poincar\'ego z odjęciem średniej
(wniosek~\ref{cor:pde10-Poincare}) pokazuje, że
$\|v\|_{V_\circ}=\|\nabla v\|_2$ jest na niej normą
równoważną normie $H^1$.
Forma $a(u,v)=\int_\Omega\nabla u\cdot\nabla v\,
\mathrm d\lambda_d$ jest ograniczona i koercywna w tej normie.
Prawa strona $F(v)$ jest liniowa i ciągła na $H^1(\Omega)$:
$|\int_\Omega fv|\le\|f\|_2\|v\|_2$,
a twierdzenie o śladzie~\ref{thm:pde10-trace} daje
$|\int_{\partial\Omega}h\operatorname{Tr}v|
\le C\|h\|_{L^2(\partial\Omega)}\|v\|_{H^1}$.
Z twierdzenia Laxa--Milgrama istnieje jednoznaczne $u\in V_\circ$
z $a(u,v)=F(v)$ dla wszystkich $v\in V_\circ$.
Dla dowolnego $v\in H^1(\Omega)$ rozkładamy
$v=(v-v_\Omega)+v_\Omega$.
Pierwszy składnik należy do $V_\circ$; dla stałej
$v_\Omega$ obie strony równania są równe zeru dzięki
warunkowi zgodności. Dostajemy równość dla każdego $v\in H^1$.

\textbf{Jednoznaczność modulo stałe.} Jeśli $u_1,u_2$
są rozwiązaniami, różnica $w=u_1-u_2$ spełnia
$\int_\Omega|\nabla w|^2\,\mathrm d\lambda_d=0$
po podstawieniu $v=w$. Wobec tego $\nabla w=0$ prawie wszędzie,
a spójność $\Omega$ implikuje, że $w$ jest stała prawie wszędzie.
\end{proof}

\section{Metoda bezpośrednia rachunku wariacyjnego}
\label{sec:pde11-direct}

Założenie, że rozwiązanie jest punktem minimum odpowiedniej energii,
pozwala badać również równania nieliniowe.
Do tego potrzebujemy precyzyjnego rozróżnienia zbieżności
silnej i słabej.

\begin{definicja}[Zbieżność słaba]
\label{def:pde11-weakconv}
W przestrzeni Hilberta $H$ piszemy $u_n\rightharpoonup u$,
jeżeli $(u_n,v)_H\to(u,v)_H$ dla każdego $v\in H$.
W przestrzeni Banacha $X$ oznacza to warunek
$F(u_n)\to F(u)$ dla każdego $F\in X^*$.
Funkcjonał $J:X\to(-\infty,+\infty]$ jest
\emph{słabo półciągły z dołu w sensie ciągów}, jeżeli
$u_n\rightharpoonup u$ implikuje
$J(u)\le\liminf_n J(u_n)$.
Jest \emph{koercywny}, jeżeli $J(u)\to+\infty$,
gdy $\|u\|_X\to\infty$.
\end{definicja}

\begin{lemat}[Słaba zwartość ograniczonych ciągów w Hilbercie]
\label{lem:pde11-weakcompact}
Każdy ciąg ograniczony w rzeczywistej przestrzeni Hilberta
$H$ zawiera podciąg słabo zbieżny.
\end{lemat}
\begin{proof}
Niech $(u_n)$ będzie ciągiem ograniczonym przez $C$.
Jego domknięta otoczka liniowa
$H_0=\overline{\operatorname{span}\{u_n:n\ge1\}}$
jest separowalną przestrzenią Hilberta.
Jeżeli $H_0$ jest nieskończeniewymiarowa,
wybieramy przeliczalną bazę ortonormalną $(e_j)$.
Każdy ciąg liczb $((u_n,e_j))_n$ jest ograniczony;
argument przekątniowy daje podciąg $(u_{n_k})$ taki, że
$(u_{n_k},e_j)\to a_j$ dla każdego $j$.
Z nierówności Bessela
$\sum_{j=1}^N|a_j|^2
=\lim_{k\to\infty}\sum_{j=1}^N|(u_{n_k},e_j)|^2\le C^2$
dla każdego $N$. Zatem $\sum_j|a_j|^2\le C^2$,
więc istnieje $u\in H_0$ o współrzędnych $(a_j)$.
Dla skończonych kombinacji wektorów $e_j$ zachodzi
$(u_{n_k},v)\to(u,v)$.
Dla ogólnego $v\in H_0$ przybliżamy go takimi kombinacjami:
$|(u_{n_k}-u,v-v_N)|\le 2C\|v-v_N\|$,
co pozwala najpierw przejść $k\to\infty$, a potem $N\to\infty$.
Dla $v\in H$ rozkład $v=P_{H_0}v+(v-P_{H_0}v)$
redukuje twierdzenie do $H_0$, bo $u_n,u\in H_0$.
W przypadku skończeniewymiarowym zwykła zwartość
ograniczonego zbioru dostarcza nawet zbieżności silnej.
\end{proof}

\begin{lemat}[Słaba półciągłość normy]
\label{lem:pde11-weaklsc}
Niech $H$ będzie rzeczywistą przestrzenią Hilberta.
Jeżeli $u_n\rightharpoonup u$ w $H$, to
$\|u\|_H\le\liminf_n\|u_n\|_H$.
Ponadto dla każdego $F\in H^*$ mamy $F(u_n)\to F(u)$.
\end{lemat}
\begin{proof}
Jeżeli $u=0$, nierówność jest oczywista.
W przeciwnym razie słaba zbieżność daje
$\|u\|_H^2=\lim_n(u_n,u)_H
\le(\liminf_n\|u_n\|_H)\|u\|_H$,
po czym dzielimy przez $\|u\|_H$.
Druga własność wynika z definicji i twierdzenia Riesza.
\end{proof}

\begin{twierdzenie}[Metoda bezpośrednia w przestrzeni Hilberta]
\label{thm:pde11-direct}
Niech $H$ będzie rzeczywistą przestrzenią Hilberta,
a $J:H\to\R$ funkcjonałem ograniczonym z dołu,
koercywnym i słabo półciągłym z dołu w sensie ciągów.
Wówczas istnieje $u\in H$ taki, że
$J(u)=\inf_{v\in H}J(v)$.
Jeżeli $J$ jest ponadto ściśle wypukły,
to minimizator jest jedyny.
\end{twierdzenie}
\begin{proof}
Niech $m=\inf_HJ\in\R$ i wybierzmy $u_n\in H$ z
$m\le J(u_n)<m+1/n$. Koercywność zapewnia
ograniczoność $(u_n)$: gdyby istniał podciąg
o normach dążących do nieskończoności,
to jego energie dążyłyby do $+\infty$, w sprzeczności z
$J(u_n)\to m$.
Na mocy lematu~\ref{lem:pde11-weakcompact}
istnieje podciąg $u_{n_k}\rightharpoonup u$.
Słaba półciągłość z dołu daje
$J(u)\le\liminf_kJ(u_{n_k})=m$.
Z definicji infimum również $m\le J(u)$;
zatem $J(u)=m$.
Jeśli istniałyby dwa różne minimizatory $u,w$,
ścisła wypukłość zapewniłaby
$J((u+w)/2)<(J(u)+J(w))/2=m$, sprzeczność.
\end{proof}

\begin{uwaga}[Dlaczego słaba półciągłość jest niezbędna?]
Sama koercywność nie wystarcza. Przykładowo na $H=\R$
zdefiniujmy $J(0)=1$, $J(x)=x^2$ dla $x\ne0$.
Funkcjonał jest ograniczony z dołu i koercywny,
lecz jego infimum wynosi $0$ i nie jest osiągane.
Nie jest też półciągły z dołu w punkcie $0$.
\end{uwaga}

\section{Energia kwadratowa i konieczność symetrii}
\label{sec:pde11-energy}

\begin{twierdzenie}[Zasada Dirichleta dla formy symetrycznej]
\label{thm:pde11-dirichlet}
Niech $H$ będzie rzeczywistą przestrzenią Hilberta,
$a:H\times H\to\R$ ograniczoną, \emph{symetryczną}
i koercywną formą dwuliniową, $F\in H^*$.
Zdefiniujmy $J:H\to\R$ wzorem
$J(v)=\frac12 a(v,v)-F(v)$.
Wtedy $J$ ma dokładnie jeden minimizator $u$, który jest
jednocześnie jedynym rozwiązaniem słabym
$a(u,v)=F(v)$ dla wszystkich $v\in H$.
Dla każdego $w\in H$ zachodzi tożsamość
\[ J(w)-J(u)=\frac12a(w-u,w-u)
\ge\frac\alpha2\|w-u\|_H^2. \]
\end{twierdzenie}
\begin{proof}
\textbf{Koercywność funkcjonału.} Z założeń otrzymujemy
\[ J(v)\ge\frac\alpha2\|v\|_H^2
 -\|F\|_{H^*}\|v\|_H\longrightarrow+\infty
 \quad\text{gdy }\|v\|_H\to\infty. \]

\textbf{Półciągłość z dołu.} Koercywna forma symetryczna
$a$ definiuje iloczyn skalarny $(v,w)_a=a(v,w)$.
Ograniczoność i koercywność dają równoważność norm:
$\sqrt\alpha\|v\|_H\le\|v\|_a\le\sqrt M\|v\|_H$.
Z twierdzenia Riesza w $(H,(\cdot,\cdot)_a)$ i
równoważności norm wynika, że zbieżność słaba w $H$
jest taka sama jak zbieżność słaba w nowym iloczynie.
Lemat~\ref{lem:pde11-weaklsc} daje więc
$\|v\|_a^2\le\liminf_n\|v_n\|_a^2$
dla $v_n\rightharpoonup v$; składnik $F(v_n)$ zbiega do $F(v)$.
Stąd $J$ jest słabo półciągły z dołu.
Twierdzenie~\ref{thm:pde11-direct} daje minimizator.

\textbf{Równanie Eulera.} Dla dowolnych $u,v\in H$
i $t\in\R$ wykorzystujemy symetrię:
$J(u+tv)=J(u)+t[a(u,v)-F(v)]+\frac{t^2}{2}a(v,v)$.
Jeżeli $u$ minimalizuje $J$, wielomian po prawej ma
minimum w $t=0$, więc jego współczynnik liniowy znika:
$a(u,v)=F(v)$. Odwrotnie, jeśli $u$ spełnia tę równość,
to po rozwinięciu $J(u+(w-u))$ otrzymujemy
$J(w)-J(u)=\frac12a(w-u,w-u)\ge0$.
Koercywność czyni nierówność ostrą dla $w\ne u$,
dając jednoznaczność. Ostatnia nierówność daje również
ilościową kontrolę błędu energetycznego.
\end{proof}

\begin{kontrprzyklad}[Dlaczego symetrii nie wolno pominąć]
\label{ex:pde11-nonsymmetric}
Na $H=\R^2$ przyjmijmy
$B=\begin{pmatrix}1&-1\\1&1\end{pmatrix}$,
$a(u,v)=(Bu)\cdot v$ oraz $F(v)=(1,0)\cdot v$.
Wtedy $a(u,u)=|u|^2$, więc $a$ jest koercywna i ograniczona.
Minimizer funkcjonału $J(u)=|u|^2/2-(1,0)\cdot u$
wynosi $u_J=(1,0)$. Tymczasem równanie
$a(u,v)=F(v)$ dla wszystkich $v$ jest równoważne
$Bu=(1,0)$ i ma rozwiązanie $u_*=\frac12(1,-1)$.
Zatem $u_J\ne u_*$. Część antysymetryczna formy znika
w $a(v,v)$, ale nie znika z samego równania słabego.
\end{kontrprzyklad}

\section[Operator p-Laplace'a]{Pierwszy problem nieliniowy: operator $p$-Laplace'a}
\label{sec:pde11-plaplace}

W części nieliniowej musimy użyć faktu, że dla
$1<p<\infty$ przestrzeń $W_0^{1,p}(\Omega)$ jest refleksywna.
Jest ona domkniętą podprzestrzenią produktu skończonej liczby
przestrzeni $L^p$, a przestrzenie $L^p$ są refleksywne w tym zakresie.
Będziemy korzystać ze standardowego twierdzenia, że każdy ciąg
ograniczony w przestrzeni refleksywnej ma słabo zbieżny podciąg
(twierdzenie Eberleina--\v Smuliana); dowód tego ogólnego wyniku
należy do analizy funkcjonalnej, nie do konstrukcji samego PDE.

\begin{twierdzenie}[Istnienie i jednoznaczność dla $p$-laplasjanu]
\label{thm:pde11-plaplace}
Niech $\Omega\subset\R^d$ będzie obszarem ograniczonym Lipschitzowskim, $1<p<\infty$, $p'=p/(p-1)$,
a $f\in L^{p'}(\Omega)$. Istnieje dokładnie jeden
$u\in W_0^{1,p}(\Omega)$ taki, że
\[
\int_\Omega |\nabla u|^{p-2}\nabla u\cdot\nabla v
\,\mathrm d\lambda_d(x)
=\int_\Omega fv\,\mathrm d\lambda_d(x)
\quad\text{dla każdego }v\in W_0^{1,p}(\Omega).
\]
W szczególności $u$ jest rozwiązaniem dystrybucyjnym
$-\operatorname{div}(|\nabla u|^{p-2}\nabla u)=f$.
Ponadto
$\|\nabla u\|_p^{p-1}\le C_P\|f\|_{p'}$.
\end{twierdzenie}
\begin{proof}
\textbf{Krok 1: energia i ograniczenie z dołu.}
Na $X=W_0^{1,p}(\Omega)$ przyjmijmy
$\|v\|_X=\|\nabla v\|_p$, co jest normą równoważną
$W^{1,p}$ dzięki nierówności Poincar\'ego.
Definiujemy
\[ J_p(v)=\frac1p\int_\Omega|\nabla v|^p
\,\mathrm d\lambda_d(x)-\int_\Omega fv
\,\mathrm d\lambda_d(x). \]
H\"older i Poincar\'e dają
$|\int_\Omega fv|\le\|f\|_{p'}\|v\|_p
\le C_P\|f\|_{p'}\|\nabla v\|_p$.
Zatem
$J_p(v)\ge\frac1p\|v\|_X^p-C_P\|f\|_{p'}\|v\|_X$,
co implikuje koercywność i ograniczenie z dołu,
ponieważ $p>1$.

\textbf{Krok 2: minimalizator.}
Niech $(v_n)$ będzie ciągiem minimalizującym.
Koercywność zapewnia ograniczoność w $X$,
a refleksywność pozwala wybrać podciąg
$v_{n_k}\rightharpoonup u$ w $X$.
Operator gradientu jest ograniczony z $X$ do $L^p(\Omega)^d$,
więc $\nabla v_{n_k}\rightharpoonup\nabla u$ w $L^p$.
Norma w przestrzeni Banacha jest słabo półciągła z dołu:
$\|z\|=\sup_{\|\ell\|_{X^*}\le1}\ell(z)$,
a przy zbieżności słabej każdy składnik supremum
przechodzi do granicy; stąd
$\|\nabla u\|_p^p\le\liminf_k\|\nabla v_{n_k}\|_p^p$.
Ponieważ funkcjonał $v\mapsto\int_\Omega fv$ jest
ciągły i liniowy na $X$, jego wartości na $v_{n_k}$
zbiegają do wartości na $u$. Zatem
$J_p(u)\le\liminf_kJ_p(v_{n_k})=\inf_XJ_p$.
Odwrotna nierówność wynika z definicji infimum,
więc $u$ minimalizuje $J_p$.

\textbf{Krok 3: wyprowadzenie słabego równania.}
Dla $u,v\in X$ rozważamy $\psi(t)=J_p(u+tv)$.
Funkcja $\xi\mapsto|\xi|^p/p$ jest klasy $C^1$ na $\R^d$
dla każdego $p>1$ i ma gradient
$|\xi|^{p-2}\xi$ (w zerze przyjmujemy wartość $0$).
W celu uzasadnienia zamiany granicy z całkowaniem
ustalmy $|t|\le1$. Twierdzenie o wartości średniej daje
\[
\left|\frac{|\nabla u+t\nabla v|^p-|\nabla u|^p}{pt}\right|
\le C_p\big(|\nabla u|^{p-1}|\nabla v|+|\nabla v|^p\big).
\]
Funkcja po prawej jest całkowalna: pierwszy składnik
należy do $L^1$ dzięki H\"olderowi z wykładnikami
$p'$ i $p$, a drugi dzięki $v\in W^{1,p}$.
Twierdzenie o zbieżności zdominowanej daje
\[ \psi'(0)=\int_\Omega |\nabla u|^{p-2}\nabla u
\cdot\nabla v\,\mathrm d\lambda_d(x)
-\int_\Omega fv\,\mathrm d\lambda_d(x). \]
Ponieważ $u$ jest minimizatorem, $\psi'(0)=0$
dla każdego $v\in X$, co daje słabe równanie.

\textbf{Krok 4: jednoznaczność.}
Funkcja $\xi\mapsto|\xi|^p/p$ jest ściśle wypukła dla
$p>1$. Jej gradient jest zatem ściśle monotoniczny:
$(|\xi|^{p-2}\xi-|\eta|^{p-2}\eta)\cdot(\xi-\eta)>0$
dla $\xi\ne\eta$. Jeśli $u_1,u_2$ są rozwiązaniami,
testujemy różnicę ich równań funkcją $v=u_1-u_2$.
Wtedy całka z nieujemnej funkcji jest równa $0$,
więc $\nabla u_1=\nabla u_2$ prawie wszędzie.
Nierówność Poincar\'ego na $X$ daje $u_1=u_2$.

\textbf{Krok 5: oszacowanie.}
W słabym równaniu podstawiamy $v=u$.
Dostajemy $\|\nabla u\|_p^p=\int_\Omega fu
\,\mathrm d\lambda_d\le C_P\|f\|_{p'}\|\nabla u\|_p$.
Dla $\|\nabla u\|_p>0$ dzielimy przez tę normę;
przypadek zerowy jest oczywisty.
\end{proof}

\begin{uwaga}[Dlaczego pominięto przypadek $p=1$?]
Przy $p=1$ energia $\int_\Omega|\nabla u|\,\mathrm d\lambda_d(x)$ nie jest ściśle wypukła,
a $W_0^{1,1}$ nie jest refleksywna. Nie można zatem
bezkrytycznie powtórzyć dowodu; pojawiają się przestrzenie
funkcji o ograniczonej wariacji ($BV$) oraz inna postać
równań Eulera--Lagrange'a.
\end{uwaga}

\section{Metoda Galerkina i pierwsze oszacowanie błędu}
\label{sec:pde11-galerkin}

\begin{twierdzenie}[Lemat C\'ei]
\label{thm:pde11-cea}
Niech $H$ będzie rzeczywistą przestrzenią Hilberta,
$a:H\times H\to\R$ formą ograniczoną ze stałą $M$
i koercywną ze stałą $\alpha>0$, a $F\in H^*$.
Niech $H_h\subset H$ będzie skończeniewymiarową
podprzestrzenią. Jeśli $u\in H$ spełnia
$a(u,v)=F(v)$ dla wszystkich $v\in H$, a $u_h\in H_h$
spełnia $a(u_h,v_h)=F(v_h)$ dla wszystkich $v_h\in H_h$,
to $u_h$ istnieje i jest jednoznaczne oraz
\[
\|u-u_h\|_H\le\frac M\alpha
\inf_{w_h\in H_h}\|u-w_h\|_H.
\]
Gdy $a$ jest symetryczna, $u_h$ jest rzutem ortogonalnym
$u$ na $H_h$ w iloczynie skalarnym $a$, a więc
$\|u-u_h\|_a=\inf_{w_h\in H_h}\|u-w_h\|_a$.
\end{twierdzenie}
\begin{proof}
Ograniczenie formy $a$ do $H_h\times H_h$ zachowuje
ograniczoność i koercywność. Z twierdzenia Laxa--Milgrama
na $H_h$ otrzymujemy dokładnie jedno $u_h$.
Odejmując równania, dostajemy \emph{ortogonalność Galerkina}:
$a(u-u_h,v_h)=0$ dla wszystkich $v_h\in H_h$.
Niech $e=u-u_h$ i $w_h\in H_h$ będzie dowolny.
Ponieważ $w_h-u_h\in H_h$, mamy
\begin{align*}
\alpha\|e\|_H^2
&\le a(e,e)=a(e,u-w_h)+a(e,w_h-u_h)\\
&=a(e,u-w_h)\le M\|e\|_H\|u-w_h\|_H.
\end{align*}
Po podzieleniu przez $\alpha\|e\|_H$ i wzięciu infimum
po $w_h$ otrzymujemy tezę; dla $e=0$ jest oczywista.
Przy symetrii $a$ jest iloczynem skalarnym.
Ortogonalność Galerkina staje się
$(u-u_h,v_h)_a=0$, więc twierdzenie Pitagorasa
daje $\|u-w_h\|_a^2=\|u-u_h\|_a^2+
\|u_h-w_h\|_a^2$ dla wszystkich $w_h\in H_h$,
z czego wynika minimalność błędu.
\end{proof}

\begin{wniosek}[Zbieżność przy zagęszczaniu przestrzeni]
\label{cor:pde11-galerkin-density}
Przy założeniach twierdzenia~\ref{thm:pde11-cea},
niech $H_n\subset H$ będą skończeniewymiarowymi
podprzestrzeniami takimi, że
$\inf_{w\in H_n}\|u-w\|_H\to0$ dla rozważanego rozwiązania $u$.
Wówczas $\|u-u_n\|_H\to0$, gdzie $u_n$ jest rozwiązaniem otrzymanym metodą Galerkina
w $H_n$.
\end{wniosek}
\begin{proof}
Prawa strona oszacowania C\'ei dąży do zera.
Nie jest potrzebne założenie, że $H_n\subset H_{n+1}$,
lecz sama gęstość \emph{sumy} podprzestrzeni bez
kontroli dla kolejnych $n$ nie wystarczyłaby.
\end{proof}

\section{Metoda elementów skończonych na odcinku}
\label{sec:pde11-fem}

Pokażemy, że teoria wariacyjna prowadzi do konkretnego układu
macierzowego. Wróćmy do $-u''=1$ na $(0,1)$ z
warunkami $u(0)=u(1)=0$. Dzielimy $[0,1]$ na
$N\ge2$ równych pododcinków długości $h=1/N$,
a węzły oznaczamy $x_i=ih$, $i=0,\ldots,N$.

Niech $V_h\subset H_0^1(0,1)$ będzie przestrzenią funkcji
ciągłych, odcinkami afinicznych, zerujących się na końcach.
Funkcje bazowe $\phi_i$, $i=1,\ldots,N-1$, spełniają
$\phi_i(x_j)=\delta_{ij}$. Przyjmując
$u_h=\sum_{i=1}^{N-1}U_i\phi_i$,
otrzymujemy układ $KU=b$, gdzie
\[
K_{ij}=\int_0^1\phi_i'\phi_j'\,\mathrm d\lambda_1,
\qquad b_i=\int_0^1\phi_i\,\mathrm d\lambda_1.
\]
Rachunek na dwóch przedziałach, na których każda
$\phi_i$ jest niezerowa, daje
$K_{ii}=2/h$, $K_{i,i+1}=K_{i+1,i}=-1/h$,
$K_{ij}=0$ dla $|i-j|>1$ oraz $b_i=h$.

\begin{stwierdzenie}[Rozwiązanie i dokładne błędy na siatce]
\label{prop:pde11-fem-errors}
Przy powyższych definicjach rozwiązaniem macierzowym jest
$U_i=u(x_i)=x_i(1-x_i)/2$, $i=1,\ldots,N-1$.
Jeśli $u_h$ jest funkcją odcinkami afiniczną o tych wartościach,
to
\[
\|u'-u_h'\|_{L^2(0,1)}=\frac h{\sqrt{12}},
\qquad
\|u-u_h\|_{L^2(0,1)}=\frac{h^2}{\sqrt{120}}.
\]
\end{stwierdzenie}
\begin{proof}
Niech $U_0=U_N=0$ oraz
$U_i=(ih)(1-ih)/2$. Ponieważ $U_i$ jest
kwadratowym wielomianem w $i$, mamy
$-U_{i-1}+2U_i-U_{i+1}=h^2$.
Po pomnożeniu przez $1/h$ dostajemy
$-U_{i-1}/h+2U_i/h-U_{i+1}/h=h=b_i$,
więc wektor $U$ rozwiązuje układ. Dodatnia określoność
macierzy $K$ wynika z
$U^TKU=\int_0^1|\sum_iU_i\phi_i'|^2\,
\mathrm d\lambda_1>0$ dla $U\ne0$;
zatem rozwiązanie jest jednoznaczne.

Na komórce $I_i=[x_i,x_{i+1}]$ wielomian $u$
ma pochodną $u'(x)=1/2-x$, a pochodna interpolantu
$u_h'$ jest stała i równa $u'(m_i)$,
gdzie $m_i=(x_i+x_{i+1})/2$.
W konsekwencji
\[
\int_{I_i}|u'-u_h'|^2\,\mathrm d\lambda_1(x)
=\int_{-h/2}^{h/2}s^2\,\mathrm d\lambda_1(s)=h^3/12.
\]
Ponieważ jest $N=1/h$ komórek, kwadrat całkowitego błędu
jest równy $h^2/12$.
Dalej na $I_i$ zachodzi dokładnie
$u(x)-u_h(x)=(x-x_i)(x_{i+1}-x)/2$.
Dlatego
\[
\int_{I_i}|u-u_h|^2\,\mathrm d\lambda_1(x)
=\frac14\int_0^h s^2(h-s)^2\,\mathrm d\lambda_1(s)
=\frac{h^5}{120}.
\]
Po zsumowaniu $N$ identycznych wkładów otrzymujemy
$\|u-u_h\|_2^2=h^4/120$, jak w tezie.
\end{proof}

\begin{center}
\begin{tikzpicture}[x=9cm,y=7cm,>=Stealth]
 \draw[->] (-.02,0)--(1.04,0) node[right] {$x$};
 \draw[->] (0,-.01)--(0,.17) node[above] {$u$};
 \draw[thick,blue,domain=0:1,samples=80] plot (\x,{\x*(1-\x)/2});
 \draw[orange!85!black,thick] (0,0)--(.25,.09375)--(.5,.125)--(.75,.09375)--(1,0);
 \foreach \x in {0,.25,.5,.75,1}{\fill[orange!85!black] (\x,{\x*(1-\x)/2}) circle[radius=.5pt];}
 \draw[blue,thick] (.13,-.033)--(.22,-.033);
 \node[black,anchor=west,font=\footnotesize] at (.23,-.033) {dokładne};
 \draw[orange!85!black,thick] (.56,-.033)--(.65,-.033);
 \node[black,anchor=west,font=\footnotesize] at (.66,-.033) {Galerkin};
\end{tikzpicture}

{\small\emph{Rysunek:} Dokładne rozwiązanie $-u''=1$ oraz
rozwiązanie Galerkina dla czterech elementów liniowych.}
\end{center}

\section{Zakres zastosowań i granice metody}
\label{sec:pde11-limits}

\begin{uwaga}[Koercywność jest warunkiem wystarczającym]
Twierdzenie Laxa--Milgrama nie obejmuje każdej formy eliptycznej.
Na przykład forma $a(u,v)=\int_\Omega\nabla u\cdot\nabla v
-\lambda\int_\Omega uv$ może przestać być koercywna
przy dużym $\lambda$. Nie oznacza to automatycznie braku
rozwiązania: zagadnienie wymaga analizy widma lub
alternatywy Fredholma. Zbadamy to w rozdziale 12.
\end{uwaga}

\begin{uwaga}[Od istnienia do regularności]
W dowodzie słabego zagadnienia Dirichleta nie założyliśmy
gładkości współczynników $A$ ani brzegu wyższej niż
Lipschitzowska. To wystarcza do istnienia rozwiązania
w $H_0^1$, ale nie gwarantuje jego regularności $H^2$.
W rozdziałach o regularności eliptycznej będziemy analizować,
które dodatkowe założenia pozwalają odzyskać drugie
pochodne i oszacowania wyższych rzędów.
\end{uwaga}

\begin{uwaga}[Co zostanie rozwinięte dalej]
Założenie $c\ge0$ i brak wyrazów pierwszego rzędu
upraszczają dowód. Dla członów dryftu potrzebne bywają
oszacowania perturbacyjne lub warunki dodatkowe.
W następnym rozdziale przejdziemy do operatorów nieograniczonych,
teorii spektralnej i półgrup; wcześniej otrzymane twierdzenia
pozwolą m.in. skonstruować odwrotność laplasjanu Dirichleta.
\end{uwaga}

% ==============================================================
% ROZDZIAŁ 12. OPERATORY NIEOGRANICZONE I METODY SPEKTRALNE
% ==============================================================
